\section{Summary}
\label{sec:2_6_summary}
In this chapter we have reviewed the bases of the two main localization methods: position fixing and dead reckoning.
As explained throughout the different sections, neither is able to offer a perfect solution but, fortunately, their strengths and weaknesses are complementary.
Position fixing methods offer absolute positioning but require the installation of infrastructure, which limits their coverage, and their accuracy will be affected by sporadic interference and complex scenarios.
The DR method is infrastructure-less, which makes it independent of the scenario. However, it only offers relative positioning and its great limitation is the unbounded accumulation of error over time.

\begin{figure}[!t]
	\centering
	\includegraphics[width=0.8\textwidth]{seamless_architecture_2}	
	\caption{Conceptual architecture of a seamless localization system}
	\label{fig:architecture}
\end{figure}
The techniques of information fusion, besides allowing to combine both methods of localization, also allow to merge a priori known information about the scenario or the nature of the movements of the object to be positioned.
It is even possible to merge the position estimates of other mobile nodes, in what is known as cooperative or network localization \cite{wymeersch_cooperative_2009}.
Therefore, it seems logical to think that the architecture of a system capable of providing localization information seamlessly should integrate all these elements, as represented in the Figure~\ref{fig:architecture}.
However, the cooperative localization, the a priori information and the position fixing methods should be considered as opportunistic, as it cannot be assumed that they will always be available.
This does not happen with DR based methods, which will work continuously, so they should be considered as the basis of a seamless localization system.

Therefore, when implementing a system with these characteristics for pedestrians, the use of wearable devices that include inertial sensors seems mandatory.
Among them, wrist-worn wearables, such as smartwatch and smartbands, stand out for their availability on the market and their convenience for users.
%\InsertFig{seamless_architecture.pdf}{fig:architecture}{Conceptual architecture of a seamless localization system}{}{1}{}